\begin{exercise}{4.1}{4}
    Man zeige den \textbf{Cosinus-Satz} $a^2 + b^2 - 2ab \operatorname{cos} \gamma = c^2$
    für jedes Dreieck einer Kongruenzebene mit positiven Seitenlängen $a, b, c$ und jeweils den entsprechenden Seiten gegenüberliegenden Winkeln $\alpha, \beta, \gamma$.
    Man zeige in denselben Notationen auch den \textbf{Sinus-Satz}
    \[
        \frac{\operatorname{sin}\alpha}{a} = \frac{\operatorname{sin}\beta}{b} = \frac{\operatorname{sin}\gamma}{c}.
    \]
    Hinweis: Man wähle eine Seite als horizontale Seite aus und berechne die Höhe
    unseres Dreiecks auf verschiedene Weisen.
\end{exercise}

\begin{minipage}{0.69\linewidth}
Seien $A, B, C$ ein Dreieck bildende Punkte auf einer Konkgruenzebene, sodass 
$|BC| = a, \quad |CA| = b, \quad |AB| = c$.
O.b.d.A. legen wir $ C = (0, 0)$ und $B = (a,0)$.
Da $|CA| = b$ und $\gamma$ bei $C$ liegt, ist
$ A = (b \operatorname{cos} \gamma, b \operatorname{sin} \gamma)$.
Sei $D$ der Lotfu\ss\space von $A$ auf die Gerade $CB$.
Dann ist $AD\perp CB$, also sind die Dreiecke $ACD$ und $ABD$ rechtwinklig und es gilt:
\[
|CD| = b \operatorname{cos} \gamma, \quad |AD| = b \operatorname{sin} \gamma, \quad \text{und} \quad |BD| = a - b \operatorname{cos} \gamma.\]
Da $ABD$ rechtwinklig ist, folgt aus dem Satz des Pythagoras:
\begin{align*}
c^2 = |AB|^2 &= |AD|^2 + |BD|^2\\
&= (b \operatorname{sin} \gamma)^2 + (a - b \operatorname{cos} \gamma)^2\\
&= b^2 \operatorname{sin}^2 \gamma + a^2 - 2ab \operatorname{cos} \gamma + b^2 cos^2 \gamma \\
&= a^2 -2ab \operatorname{cos} \gamma + b^2(\operatorname{sin}^2 \gamma + \operatorname{cos}^2 \gamma)\\
&= a^2 + b^2 -2ab \operatorname{cos} \gamma
\end{align*}
Sei nun $h=|CE|$ die Höhe auf der Seite $c$ mit $E$ als Lotfu\ss \space von $C$ auf $c$. Dann kann $h$ im Dreieck $ACE$ und $BCE$ berechnet werden. Also gilt:
\[h=b\operatorname{sin}\alpha \qquad \text{und} \qquad h=a\operatorname{sin}\beta.\]
\end{minipage}
\begin{minipage}{0.3\linewidth}
    \begin{center}
        \begin{tikzpicture}[
        scale=1,
        point/.style={circle, fill, inner sep=1.5pt}
    ]

    % Punkte
    \coordinate (C) at (0,0);
    \coordinate (B) at (3,0);
    \coordinate (A) at (1.5,2.25);
    \coordinate (D) at (1.5,0);
    \coordinate (E) at (2.25,1.125);

    % Dreieck
    \draw[thick]
        (B) -- node[below] {$a$}
        (C) -- node[left, sloped, above] {$b$}
        (A) -- node[right, sloped, above] {$c$}
        cycle;

    % Lotfuß
    \draw[thick, red!40] (A) -- node[midway, right, text=red!40] {$h'$} (D);
    \draw[thick, green!40] (C) -- node[midway, above, text=green] {$h$} (E);

    % Punkte markieren
    \node[point, label=below left:$C$] at (C) {};
    \node[point, label=below right:$B$] at (B) {};
    \node[point, label=above:$A$] at (A) {};
    \fill[red!40] (D) circle (1.5pt) node[above right, text=red!40] {$D$};
    \fill[green] (E) circle (1.5pt) node[below, text=green] {$E$};

    % Winkel gamma bei C
    \pic[
        draw,
        angle radius=3mm,
        angle eccentricity=1.4,
        pic text={$\gamma$}
    ] {angle=B--C--A};

    % Winkel beta bei B
    \pic[
        draw,
        angle radius=3mm,
        angle eccentricity=1.4,
        pic text={$\beta$}
    ] {angle=A--B--C};

    % Winkel alpha bei A
    \pic[
        draw,
        angle radius=3mm,
        angle eccentricity=1.4,
        pic text={$\alpha$}
    ] {angle=C--A--B};

    \pic[
        draw,
        angle radius=2mm,
        angle eccentricity=0.5,
        pic text={\fontsize{3}{3}\selectfont$\bullet$}
    ] {angle=A--D--C};

    \pic[
        draw,
        angle radius=2mm,
        angle eccentricity=0.5,
        pic text={\fontsize{3}{3}\selectfont$\bullet$}
    ] {angle=A--E--C};
            
    \end{tikzpicture}
    %\begin{tikzpicture}[
    scale=1,
    point/.style={circle, fill, inner sep=1.5pt}
    ]

    % Punkte
    \coordinate (C) at (0,0);
    \coordinate (B) at (3,0);
    \coordinate (A) at (3,2);

    % Dreieck
    \draw[thick]
        (B) -- node[below] {$a$}
        (C) -- node[left, sloped, above] {$c$}
        (A) -- node[right, sloped, above] {$b$}
        cycle;

    % Winkel gamma bei C
    \pic[
        draw,
        angle radius=5mm,
        angle eccentricity=1.4,
        pic text={$\alpha$}
    ] {angle=B--C--A};

    % Winkel beta bei B
    \pic[
        draw,
        angle radius=3mm,
        angle eccentricity=0.5,
        pic text={\fontsize{3}{3}\selectfont$\bullet$}
    ] {angle=A--B--C};

    % Winkel alpha bei A
    \pic[
        draw,
        angle radius=4mm,
        angle eccentricity=1.7,
        pic text={$\beta$}
    ] {angle=C--A--B};

    \end{tikzpicture}
    \fcolorbox{gray}{white}{
        $\displaystyle
        \begin{aligned}
            \sin \alpha &= \frac{Gegenkathete}{Hypotenuse}\\
            \cos \alpha &= \frac{Ankathete}{Hypotenuse}\\
            \tan \alpha &= \frac{\sin \alpha}{\cos \alpha} =\frac{GK}{AK}\\
            \sin^2 &\alpha + \cos^2 \alpha = 1
        \end{aligned}
        $
    }
    \end{center}
\end{minipage}
Also ist
\[ \frac{\operatorname{sin}\alpha}{a} = \frac{\operatorname{sin}\beta}{b}.\]
Betrachtet man analog nun die Höhe $h'$ auf der Seite $a$, dann kann $h'$ im Dreieck $ACD$ und $ABD$ berechnet werden. Also gilt:
\[h' = c \operatorname{sin}\beta = b \operatorname{sin}\gamma\]
Also folgt
\[ \frac{\operatorname{sin}\beta}{b} = \frac{\operatorname{sin}\gamma}{c}.\]

\begin{exercise}{4.2}{3}
    Man zeige: Je zwei angeordnete Dreiecke
    \((A,B,C)\) und \((A',B',C')\)
    einer Kongruenzebene mit zwei gleichen Seitenlängen
    \(a=a'\) sowie \(b=b'\)
    und gleichem eingeschlossenen Winkel
    \(
        \gamma=\gamma'
    \)
    sind angeordnet kongruent.
    Gilt statt \(\gamma=\gamma'\) eine der Gleichheiten
    \(
        \alpha=\alpha'\) oder \(\beta=\beta',
    \)
    so brauchen unsere angeordneten Dreiecke nicht kongruent zu sein.
\end{exercise}

\[
    \boxed{\text{Dreieckswinkel liegen im Intervall } (0, \pi) \quad \Longrightarrow \quad \text{Cos ist eindeutig, Sin nicht!}}
\]
Nach dem Kosinussatz gilt
\[
    c^2=a^2+b^2-2ab\cos\gamma \qquad \text{und} \qquad
    (c')^2=(a')^2+(b')^2-2a'b'\cos\gamma'.
\]
Aus \(a=a', b=b', \gamma=\gamma'\)
folgt daher \(c^2=(c')^2.\)
Da Seitenlängen positiv sind, gilt \(c=c'\). Somit stimmen alle drei
entsprechenden Seitenlängen überein, und nach dem Kongruenzsatz SSS sind
die beiden angeordneten Dreiecke angeordnet kongruent.

Dass die Behauptung bei \(\alpha=\alpha'\) im Allgemeinen nicht gilt,
zeigt der Sinussatz. Wähle
\[
    a=1,\qquad b=\sqrt3,\qquad \alpha=30^\circ.
\]
Dann gilt
\[
    \sin\beta
    =
    \frac{b}{a}\sin\alpha
    =
    \frac{\sqrt3}{2}.
\]
Daher sind sowohl \(\beta_1=60^\circ\) als auch \(\beta_2=120^\circ\) möglich.
Die beiden Dreiecke besitzen also dieselben Werte für \(a\), \(b\) und
\(\alpha\), sind aber wegen \(\beta_1\neq \beta_2\) nicht kongruent.

Der Fall \(\beta=\beta'\) folgt analog durch Vertauschen der Rollen von
\(a,\alpha\) und \(b,\beta\).

\[
    \boxed{\text{SSS lässt sich mithilfe vom Cosinussatz und SWS (diesem Satz) genau so beweisen!}}
\]

\begin{exercise}{4.3}{3}
    (\textbf{Winkelsumme in konvexen Vielecken}). Man zeige: Seien in einer
    Kongruenzebene $E$ paarweise verschiedene Punkte $p_1, ... , p_n$ mit $n \geq 3$ und
    in $\mathbb Z/n\mathbb Z$ verstandenen Indizes gegeben derart, dass die halboffenen Segmente
    $[p_{i-1}, p_i)$ paarweise disjunkt sind und $(p_{i+1} - p_i, p_{i-1} - p_i)$ für alle $i$ linear unabhängig sind und als angeordnete Basen alle zu derselben Orientierung von $\vec{E}$
    gehören. So gilt in der Überlagerung $\tilde{\mathbb W}$ der Winkelgruppe
    \[
        \sum_{i=1}^{n}\angle (p_{i+1}-p_i, p_{i-1}-p_i) = (n-2) 180^\circ
    \]
    Man zeige durch ein Beispiel, dass die Orientierungsbedingung notwendig ist.
\end{exercise}

Sei $\phi_i = \angle (p_{i+1}-p_i, p_{i-1}-p_i)$ der Innenwinkel und $\varphi_i$ der Au\ss enwinkel bei $p_i$.
Es gilt:
\[
    \varphi_i = 180^\circ - \phi_i \quad \text{,also auch} \quad \phi_i = 180^\circ - \varphi_i \quad \text{für alle } i = 1, ..., n.
\]
Wegen der Orientierungsbedingung erfolgen alle Richtungsänderungen beim
Umlaufen des Vielecks im gleichen Drehsinn. Nach einem vollständigen Umlauf
hat sich die Richtung insgesamt um \(360^\circ\) gedreht.
Wir wissen daher
\[
    \sum_{i=1}^{n} \varphi_i = 360^\circ.
\]
Also gilt:
\[
    \sum_{i=1}^{n}\phi_i = \sum_{i=1}^{n}( 180^\circ - \varphi_i) = \sum_{i=1}^{n}180^\circ - \sum_{i=1}^{n} \varphi_i = n \cdot 180^\circ - 360^\circ = (n-2) 180^\circ.
\]

\begin{minipage}{0.69\linewidth}
    Beweis, dass die Orientierung wichtig ist:

    Wähle $E = \mathbb R^2$ mit
    \[
        p_1 = (0,1), p_2 = (-1, -1), p_3 = (0, 0) \quad \text{und} \quad p_4 = (1, -1).
    \]
    Somit sind alle Vorraussetzungen au\ss er die Orientierung erfüllt.
    Nach $\angle (p_{i+1}-p_i, p_{i-1}-p_i)$ erhalten wir die Winkel \(53.1^\circ + 18.4^\circ + 18.4^\circ + 90^\circ = 180^\circ \neq (n-2)180^\circ = 360^\circ\)
\end{minipage}
\begin{minipage}{0.3\linewidth}
    \begin{center}
        \begin{tikzpicture}[
        scale=1,
        point/.style={circle, fill, inner sep=1.5pt}
    ]

    % Punkte
    \coordinate (P1) at (0,1);
    \coordinate (P2) at (-1,-1);
    \coordinate (P3) at (0,0);
    \coordinate (P4) at (1,-1);

    % Dreieck
    \draw[thick]
        (P1) -- 
        (P2) --
        (P3) --
        (P4) --
        cycle;

    % Punkte markieren
    \node[point, label=above:$p_1$] at (P1) {};
    \node[point, label=below left:$p_2$] at (P2) {};
    \node[point, label=below:$p_3$] at (P3) {};
    \node[point, label=below right:$p_4$] at (P4) {};

    % Winkel
    \pic[
        draw,
        angle radius=4mm,
        angle eccentricity=1.4,
    ] {angle=P3--P2--P1};
    \pic[
        draw,
        angle radius=2mm,
        angle eccentricity=1.4,
    ] {angle=P2--P3--P4};
    \pic[
        draw,
        angle radius=3mm,
        angle eccentricity=1.4,
    ] {angle=P2--P1--P4};
    \pic[
        draw,
        angle radius=4mm,
        angle eccentricity=1.4,
    ] {angle=P1--P4--P3};
            
    \end{tikzpicture}
\end{center}
\end{minipage}


\begin{exercise}{4.4}{3}
    Man zeige die Formel für den Flächeninhalt eines Dreiecks einer Kongruenzebene in der Gestalt
    \[
        \sin(\alpha)bc = \pm F(B-A,C-A)
    \]    
    mit unserem orientierten Flächeninhalt auf der rechten Seite.
\end{exercise}

Schreibe $v:=B-A$ und $w:= C-A$, also $|v| = c$ und $|w| = b$. Zerlege $v$ in $v = v^+ + v^-$, wobei $v^+$ ein skalares Vielfaches von $w$ und $v^-$ senkrecht auf $w$ ist.
Nach dem Satz des Pythagoras haben wir also ein rechtwinkliges Dreieck $|v|^2 = |v^+|^2 + |v^-|^2$ konstruiert und wissen $|v^-| = |v| \cdot \sin(\alpha) = c \cdot \sin(\alpha)$.
Dann gilt:
\[
    F(v,w)=F(v^++v^-,w) = F(v^+, w) + F(v^-, w) \overset{\star}{=} F(v^-, w) \qquad ({\scriptstyle\star}: v^+ =  \lambda w).
\]
Weiter gilt:
\[
    |F(v^-,w)| = \|v^-\|\cdot \|w\|= c \sin(\alpha) \cdot b = \sin(\alpha)bc.
\]
Es folgt die Behauptung.